\section{Literatura relacionada}\label{sec:literature}

Este artículo se vincula con tres líneas de investigación: la teoría y la evidencia sobre
los efectos del salario mínimo en el empleo, la evidencia para países en desarrollo con
amplios sectores informales y la literatura específica sobre el Perú.

\paragraph{Teoría y evidencia canónica.} El modelo competitivo predice que un piso salarial
vinculante reduce el empleo de los trabajadores de baja productividad, con un efecto cuyo
tamaño depende de las elasticidades de la demanda de trabajo y de la fracción de
trabajadores alcanzados por el piso \citep{brown_1999}. Los modelos en los que los
empleadores tienen poder para fijar salarios generan predicciones distintas: bajo monopsonio,
un salario mínimo moderado puede elevar el empleo y, por lo general, comprime la cola
inferior de la distribución salarial \citep{manning_2003}. La literatura empírica ha pasado
de presumir efectos de desempleo considerables a un consenso de efectos pequeños en el
margen estudiado. El estudio de caso de \citet{card_krueger_1994} y la síntesis más amplia de
\citet{card_krueger_1995} encontraron poca evidencia de pérdida de empleos por la reforma de
Nueva Jersey, y tanto el diseño de discontinuidad fronteriza de
\citet{dube_lester_reich_2010} como el enfoque de acumulación (\emph{bunching}) de
\citet{cengiz_2019} hallan igualmente que los aumentos recientes en Estados Unidos elevaron
los salarios en la parte baja de la distribución con efectos limitados sobre el empleo;
\citet{manning_2021} interpreta la evidencia acumulada como un efecto sobre el empleo
``esquivo'', mientras que \citet{neumark_wascher_2007} ofrecen la lectura escéptica.
\citet{harasztosi_lindner_2019} muestran que un gran aumento en Hungría se absorbió
principalmente mediante precios y no mediante empleo. Lo más cercano a mi mecanismo es
\citet{dustmann_2022}, quienes muestran que el salario mínimo alemán operó mediante una
\emph{reasignación} que movió trabajadores entre establecimientos, y no mediante pérdidas de
empleo, y \citet{engbom_moser_2022}, quienes documentan cómo el aumento del mínimo brasileño
comprimió la distribución salarial formal; \citet{derenoncourt_montialoux_2021} rastrean las
consecuencias distributivas de las ampliaciones de cobertura en Estados Unidos. Un tema
recurrente, sobre el que construyo directamente, es que los efectos deben medirse allí donde
el mínimo tiene incidencia (\emph{bite}). \citet{card_1992} usa la variación regional en la
fracción de trabajadores afectados, \citet{clemens_wither_2019} siguen a trabajadores cuyos
salarios se encontraban por debajo de un piso federal creciente, \citet{lee_1999} usa la
brecha entre el mínimo y el salario mediano, y \citet{machin_manning_rahman_2003} aprovechan
un sector de bajos salarios en el que la introducción del mínimo en el Reino Unido fue
especialmente vinculante. Mi diseño basado en la educación aplica la misma lógica a una
única reforma nacional.

\paragraph{Países en desarrollo e informalidad.} En economías con amplios sectores
informales, el salario mínimo interactúa con la frontera entre el trabajo cubierto y el no
cubierto, y aparecen repetidamente dos fuerzas opuestas. Un efecto faro puede elevar los
salarios incluso en el sector informal, donde el mínimo funciona como norma o referencia
\citep{lemos_2009, bosch_manacorda_2010, maloney_nunez_2004}, mientras que un efecto de
reasignación puede expulsar a trabajadores del empleo formal hacia arreglos informales o de
trabajo independiente \citep{ham_2018_honduras}; \citet{ulyssea_2018} ofrece el marco de
equilibrio moderno en el que la regulación desplaza a empresas y trabajadores a través de la
frontera de la formalidad. La evidencia latinoamericana abarca todo este rango. Los primeros
trabajos sobre México y Colombia encontraron efectos de desempleo donde el piso era alto en
relación con los salarios \citep{bell_1997}, y \citet{gindling_terrell_2007} muestran para
Costa Rica que los pisos legales mueven los salarios solo en el sector cubierto. Los estudios
sobre Brasil hallan compresión salarial con efectos pequeños sobre el empleo y evidencia de
efectos de derrame hacia el sector informal \citep{lemos_2009, lemos_2004};
\citet{jales_2018} estima, con un enfoque de densidades, que el mínimo brasileño traslada a
una proporción considerable de trabajadores a la informalidad. Los trabajos sobre Ecuador
encuentran efectos salariales positivos concentrados entre los trabajadores de bajos
ingresos \citep{wong_2019_ecuador}, y \citet{arango_2022_colombia} reúnen la evidencia
colombiana reciente. El caso hondureño es el contraejemplo más claro a la visión del efecto
faro: \citet{ham_2018_honduras} documenta que un salario mínimo alto redujo el empleo
cubierto y desplazó a los trabajadores hacia la informalidad, sin reducir la pobreza. Para
Uruguay, \citet{katzkowicz_2021_uruguay} usan un enfoque de discontinuidad en la densidad en
el sector del trabajo doméstico para mostrar que los efectos se concentran entre los menos
calificados. El hilo metodológico común es el uso de medidas de exposición o de incidencia,
ya sean índices de Kaitz geográficos o proporciones de fracción afectada, y una
contabilización cuidadosa del margen informal; adopto ambos elementos.

\paragraph{Perú.} La literatura nacional es escasa y, hasta hace poco, desactualizada; sin
embargo, el margen que enfatizo ha comenzado a recibir atención: \citet{corcuera_2024}
estudia la respuesta del trabajo independiente informal a los aumentos del salario mínimo
peruano con datos anteriores. Los primeros estudios de las reformas de inicios de la década
de 2000 hallan efectos pequeños y a menudo estadísticamente débiles.
\citet{jaramillo_lopez_2006} analizan el aumento de 2003 con paneles trimestrales cortos para
Lima Metropolitana y concluyen que la reforma fue en gran medida inocua, con efectos
negativos sobre la permanencia en el empleo concentrados cerca del piso y evidencia de
desplazamiento del sector formal al informal. \citet{cespedes_2005} estima una elasticidad
del empleo débilmente significativa, de alrededor de $-0.13$, y muestra que la RMV causa en
el sentido de Granger a los salarios formales, una formulación temprana de la visión del
efecto faro para el Perú. El estudio moderno más creíble es \citet{jimenez_2023}, que evalúa
la reforma de 2016 con un diseño de diferencias en diferencias y una dosis de tratamiento a
nivel departamental, definida como la proporción previa a la reforma de trabajadores
formales que ganaban entre el mínimo anterior y el nuevo. Ese estudio no encuentra efectos de
desempleo, halla un aumento de la formalidad de 2.5 a 2.8 puntos porcentuales y ganancias
salariales de 4.1 a 5.4 por ciento, e interpreta este patrón a través del poder de mercado de
los empleadores, del tipo que \citet{amodio_deroux_2021} documentan directamente en plantas
manufactureras colombianas. Mi artículo se diferencia de este precedente en tres aspectos que
hago explícitos. Estudio una reforma posterior y el período posterior a la pandemia; defino
la exposición a nivel de trabajador mediante la educación y no a nivel departamental; y
encuentro un patrón de ajuste distinto, sin ganancia salarial y con una caída, y no un
aumento, de la formalidad. Vuelvo a la conciliación de estos resultados en la
Sección~\ref{sec:discussion}. Como enfatiza \citet{jimenez_2023}, la restricción decisiva para
cualquier diseño peruano es la ausencia de variación subnacional en el mínimo nominal, lo que
obliga al investigador a construir una exposición diferenciada; mi contribución es hacerlo a
través de la dimensión educativa y rastrear sus consecuencias en los márgenes formal e
informal.

\paragraph{Econometría de diferencias en diferencias.} En lo metodológico me apoyo en la
literatura moderna de diferencias en diferencias. Como mi reforma ocurre en una sola fecha,
las preocupaciones recientes sobre ponderaciones negativas y comparaciones prohibidas bajo
adopción escalonada \citep{goodmanbacon_2021, dechaisemartin_2020, sun_abraham_2021,
callaway_santanna_2021} no aplican: con una fecha de tratamiento común, el estimador de
efectos fijos bidireccionales y su análogo de estudio de eventos son insesgados para un
promedio convexo de efectos por grupo y período \citep{roth_2023_trending,
baker_larcker_wang_2022}. Aun así, uso el estimador doblemente robusto de
\citet{santanna_zhao_2020} como pieza de base, y tomo en serio la amenaza a la
identificación, es decir, las tendencias diferenciadas por calificación, mediante el
análisis de sensibilidad de \citet{rambachan_roth_2023}. Para la inferencia con pocos
conglomerados sigo a \citet{bertrand_2004}, \citet{cameron_2008_bootstrap} y
\citet{pustejovsky_tipton_2018}; la selección en la muestra salarial se acota con el enfoque
de recorte de \citet{lee_2009}. Los efectos distributivos se estiman con el método de
cuantiles incondicionales de \citet{firpo_fortin_lemieux_2009}.
